\section{Trabajo futuro}

La mejora más directa es reducir el efecto del cono del sensor. El VL53L1X, de
la misma familia, permite elegir una región de interés más pequeña dentro de su
campo de visión, lo que haría los bordes más nítidos. Con el sensor actual, una
alternativa por software es la deconvolución: si se conoce la forma del cono,
la relación entre la superficie real y las mediciones es un sistema lineal que
se puede invertir de manera aproximada por mínimos cuadrados con
regularización.

Para obtener un modelo cerrado del objeto y no solo su vista desde arriba, se
puede escanear desde varios lados y unir las nubes de puntos. Cada escaneo
queda en su propio sistema de coordenadas, y la rotación y la traslación que
los alinean se pueden estimar con el algoritmo ICP, en el que la rotación
óptima de cada paso se obtiene con la descomposición en valores singulares
\citep{besl1992,arun1987}.

Por último, la posición de la caja podría medirse en lugar de depender de la
cuadrícula, por ejemplo con dos sensores de distancia más apuntando a las
paredes de una guía, lo que permitiría mover la caja libremente y tomar
mediciones de forma continua.
